% Collection of random stuff we might want later.


\klcomment{feels like appendix content}
While estimating source trustworthiness has seen more attention in fake news identification \cite{baly-etal-2018-predicting}, some have shown it to also be useful in fact-checking, especially in web and social media domains \cite{Popat2017WhereTT,Nguyen2018AnIJ}.  Like \fever and \snopes, \ours does not explicitly handle this as our sources of documents are peer-reviewed academic journals.




%%%%%%%%%%%%%%%%%%%%%%%%%%%%%%%%%%%%%%%%%%%%%%%%%%%%%%%%%%%%%%%%%%%%%%%%%%%%%%%%

% \paragraph{Statistics}

% \kyle{Stopped here;  Need more info about train-dev-test splits for full corpus not just claim-abstract pairs} \dw{The corpus isn't split.}

% We divide the collection of claim-abstract pairs into standard train-dev-test splits.  
% The label distribution is shown in Table \ref{tbl:claim_counts}. 
% % \footnote{The test set is balanced, but the dev set is not; due to the smaller size of our dataset, we want to make sure our models have as many \refutes labels to train on as possible.} 

% Some of these medical journals feature \emph{structured} abstracts, which are organized into sections focusing on background, methods, results, and conclusions. The majority of abstracts are \emph{unstructured}. As shown in Table \ref{tbl:dataset_overview}, structured abstracts are generally longer and more likely to contain multiple rationales. The majority of claims have evidence in just a single document. 

% \begin{table}[t]
  \setlength{\tabcolsep}{.25em}
  \footnotesize
  \centering

  \begin{tabular}{L{0.12\linewidth} C{0.19\linewidth} C{0.19\linewidth} C{0.19\linewidth} C{0.15\linewidth}}
    \toprule
    Fold &  \supports &  \textsc{NEI} &  \refutes &   All \\
    \midrule
    Train &       332 &              304 &      173 &   809 \\
    Dev   &       124 &              112 &       64 &   300 \\
    Test  &       100 &              100 &      100 &   300 \\
    \midrule
    All   &       556 &              516 &      337 &  1409 \\
    \bottomrule
  \end{tabular}

  \caption{Distribution of claim / cited abstract labels by train-dev-test split in \ours. \textsc{NEI} stands for \notenough.}
  \label{tbl:claim_counts}
  
\end{table}


% For detailed dataset statistics, see Appendix \ref{sec:corpus_stats}.


%%%%%%%%%%%%%%%%%%%%%%%%%%%%%%%%%%%%%%%%%%%%%%%%%%%%%%%%%%%%%%%%%%%%%%%%%%%%%%%%

% Label agreement

Treating the first annotator's label as gold and the second as a prediction, the label classification accuracy is \labelhumanacc. \kyle{what does this tell us that Cohen's doesn't?}



%%%%%%%%%%%%%%%%%%%%%%%%%%%%%%%%%%%%%%%%%%%%%%%%%%%%%%%%%%%%%%%%%%%%%%%%%%%%%%%%

% BioMedSumm

This represents a substantial improvement over the BioMedSumm citation contextualization dataset, for which the inter-annotator agreement in selection rationale sentences is 0.27 Fleiss' $\kappa$\footnote{BioMedSumm rationales are annotated at the span level; we convert to sentence-level }